\section{Conclusion and Future Work}
\label{submission:sec:conclusion}

In this paper, we challenged the foundational assumption of \emph{spatial homogeneity} in stateful dynamic model merging. We introduced \textbf{Layer-Decoupled Stateful Kinetics (LDS-Kinetics)}, a learning-theoretic framework that partitions dynamic ensembling layers into depth-wise decoupled blocks. Each block maintains an independent concentration state that evolves dynamically according to local workload similarity and block-specific parameters. To counter the threat of transductive overfitting in highly decoupled spaces, we derived a unified PAC-Bayesian complexity penalty across all blocks simultaneously.

Guided by our empiricist philosophy, we evaluated LDS-Kinetics across extensive multi-dimensional parameter sweeps inside a 14-layer analytical coordinate sandbox. Our results show that LDS-Kinetics achieves a superior accuracy-jitter Pareto frontier compared to prior state-of-the-art stateful routing methods. Crucially, our empirical analysis deconstructed stateful ensembling depth dynamics for the first time, demonstrating that neural networks naturally organize temporal scales along depth: early blocks maintain rapid, dynamic tempos to perform spatial manifold alignment, whereas late blocks develop low decay (high inertia) to act as stable low-pass filters for decision logits.

Future work will focus on three key directions. First, as our highest-priority immediate next step, we will validate LDS-Kinetics on physical pre-trained models (such as LLaMA-3, Mistral, and Vision Transformers) serving real-world sequential task streams, confirming that the depth-dependent "tempo-gradient" generalizes to physical activation manifolds. As a step towards this, we have successfully implemented and verified a complete physical integration proof-of-concept in PyTorch (\texttt{test\_physical\_poc.py}) demonstrating dynamic weight-blended LoRA expert adapters with perfect shape and simplex constraint satisfaction. Second, we plan to execute comprehensive systems-level GPU latency benchmarks, evaluating the empirical throughput and CUDA overhead of batched tensor state updates and fused Triton-based kernels to guarantee negligible serving latency under high-concurrency production workloads. Finally, we will explore adaptive block partitioning schemes that dynamically learn optimal layer-wise groupings online based on real-time activation statistics rather than static pre-partitions.
